\begin{problem}[Affine anti-equivalence]\label{prob:vi19-p04}
Prove the natural bijection
\[
\operatorname{Hom}_{\mathrm{AffSch}}(\Spec B,\Spec A)
\cong
\operatorname{Hom}_{\mathrm{Ring}}(A,B),
\]
and show that composition corresponds to reversed composition of ring maps.
\end{problem}

\ifdefined\FullProblemDossiers
\paragraph{Why this problem matters.}
This is the central computational principle for affine geometry.

\paragraph{Useful ideas before starting.}
Apply VI.19.P1 with $X=\Spec B$ and use $\Gamma(\Spec B,\OO)=B$.

\paragraph{Guided hints.}
\begin{enumerate}
\item Specialize the universal property.
\item Use VI/18 global recovery.
\item Track composition on global sections.
\end{enumerate}

\begin{solution}
% VI19 full-solutions refinement v1
Apply the universal property VI.19.P1 to the locally ringed space \(X=\Spec B\).  It gives a natural bijection
\[
\operatorname{Hom}_{\mathrm{LRS}}(\Spec B,\Spec A)
\cong
\operatorname{Hom}_{\mathrm{Ring}}(A,\Gamma(\Spec B,\OO)).
\]
By affine global recovery from VI/18,
\[
\Gamma(\Spec B,\OO)\cong B,
\]
canonically.  Since morphisms between affine schemes are precisely morphisms of locally ringed spaces between these objects, we obtain
\[
\boxed{
\operatorname{Hom}_{\mathrm{AffSch}}(\Spec B,\Spec A)
\cong
\operatorname{Hom}_{\mathrm{Ring}}(A,B).}
\]

The correspondence sends a ring map \(\varphi:A\to B\) to \(\Spec(\varphi):\Spec B\to\Spec A\), and sends an affine morphism to its pullback on global sections.  VI.19.P1 proves these constructions are inverse, so the bijection is not merely a counting statement.

Now let \(A\xrightarrow{\varphi}B\xrightarrow{\psi}C\).  The corresponding geometric arrows are \(\Spec B\to\Spec A\) and \(\Spec C\to\Spec B\).  Their composite is \(\Spec(\psi\circ\varphi)\) by functoriality.  Hence geometric composition corresponds to ordinary ring composition read in the reverse geometric direction.  This is exactly the assertion that \(\mathbf{AffSch}\simeq(\mathbf{CRing})^{\mathrm{op}}\).
\end{solution}

\paragraph{What happened mathematically.}
The category of affine schemes is the opposite of the category of commutative rings.

\paragraph{Extensions.}
Derive the affine-scheme isomorphism criterion.

\paragraph{Provenance.}
\textbf{Canonical derivative of AG3 Problem 12 / VI.19.P1; not an additional legacy migration.}
\else
\paragraph{Provenance.} \textbf{Canonical derivative of AG3 Problem 12 / VI.19.P1; not an additional legacy migration.}
\fi
